The axis of Section~2 is only useful if it can be applied before the work it is
meant to inform. This section gives the procedure that does so, and its defining
property is cost: it reads source and nothing else. No system is built, nothing is
executed, and no test campaign is run. The output is an upper bound on an oracle's
reach, obtained while the oracle is still a design.

\subsection{Inputs}

The procedure takes two inputs.

The first is the oracle \emph{design} $O$, characterized not by its implementation
but by its \textbf{consumption set}: the artifacts $O$ will read in order to form
its judgement. Naming this set requires no code, which is what makes the procedure
a pre-build one. The set must be named at the granularity of individual artifacts
or fields---``the profile of executed basic blocks'', ``the range fields of the
tracked abstract state''---and not as ``the system's output'', because the answer
differs per artifact. It must also be named relative to the \emph{component} under
test rather than the whole system, for the reason given in Section~2: an oracle may
share most of a system with the code it judges and still be evidentially
independent of the part that matters.

The second is a defect population $P$ drawn from the system under test, and it
carries a precondition.

\subsection{Precondition on the population}

\textbf{$P$ must be drawn independently of any oracle's findings.} A population
assembled from what some tool reported measures the sighted side of that tool and
cannot bound the blind side of any tool, including itself. The natural constructions
that satisfy the precondition are exhaustive over some independently chosen frame:
every fixed defect in a release window, or every fix touching a named component in
the project's history. The natural construction that violates it is the appealing
one---the bug list a technique published---because it is already assembled,
already curated, and already classified.

This is a precondition and not a caution: if $P$ is drawn from a tool's findings,
the procedure does not yield a weaker bound, it yields no bound at all. We state it
formally because it does not transfer on its own. Having applied it correctly to the
system in Section~5---partitioning the historical population by what an oracle
\emph{could} observe, not by what it had found---we nonetheless proposed, when
moving to a second system, to classify the defect list that a published technique
had itself reported. The asymmetry had been learned in one context and did not
survive the move to another, which is the argument for writing the procedure down
rather than treating it as judgement.

A second requirement on $P$ is weaker but must be reported: the criterion selecting
$P$ should not correlate with visibility. A frame such as ``defects whose fix
message quotes internal state'' is convenient and biased, because the property that
makes a defect easy to read may also be the property that made it findable.

\subsection{The step}

For each defect $d \in P$, ask one question:

\begin{quote}
Does $d$ corrupt an artifact in $O$'s consumption set?
\end{quote}

\noindent Three answers are permitted.

\begin{itemize}
\item \textbf{Yes} --- $O$ consumes the corruption and reproduces it. $d$ is in
      $O$'s blind class, structurally and regardless of $O$'s predicate.
\item \textbf{No} --- $d$ corrupts nothing $O$ consumes. This answer covers two
      situations that must be kept apart. If $d$ lies in an artifact $O$ never reads
      at all, $O$ is blind to it as surely as in the first case, for a different
      reason. If $d$ lies in something $O$ does examine, $O$ \emph{may} observe it,
      subject to its predicate covering the class---a necessary condition, not a
      sufficient one. The reproducibility measurement of Section~4.5 found that
      independent readers split exactly on this seam, which is why it is named here.
\item \textbf{Undecided} --- the fix does not settle the question from source.
\end{itemize}

Two disciplines make the step reliable, and both were earned rather than assumed.
First, \emph{decide from the source of the fix, never from its title or summary}.
In an earlier partition of forty fixes, four carried titles that named a
component the diff did not touch; a classification taken from titles would have been
wrong in one case in ten. Second, \emph{Undecided is a bucket, not a default}. A
procedure that silently resolves ambiguity produces a bound it cannot support, and
a count of undecided cases is part of the result rather than an embarrassment in it.

\subsection{Output, and what it is not}

The partition yields $\mathrm{reach}(O) \le |\{d : \text{No}\}| / |P|$, an
\emph{upper} bound, and the slack in it has a name: predicate coverage. A defect the
oracle does not consume may still lie in an arm the oracle never implemented, and
$|\{\text{No}\}|$ counts it anyway. In Section~5 the bound is $1/57$ and the
realized reach is $0/57$; the difference is one such arm. This slack is not
irreducible before a build, because which arms an oracle implements is a property
of the oracle's source, readable the same way the partition is. Being outside the
blind class is necessary for observation and
not sufficient: the oracle's predicate must also cover the class. In Section~5 the
bound is one defect in fifty-seven, and the realized reach is zero, because that one
lies in a part of the predicate the oracle did not implement. A procedure that
promised reach rather than bounding it would have promised one defect and delivered
none on its first application.

\subsection{Why the cost matters}

The procedure's value is not that it is rigorous but that it is cheap enough to run
before committing. Calibrating a single oracle against a single real defect, in the
manner of Section~4, costs a build of the system at the defect and another at its
fix; in our case study that is tens of minutes of machine time per pair and a
configuration that must be held byte-identical across the two. The procedure costs
opening one commit per defect and reading the function it changes.

The asymmetry is the point: run the procedure over the whole population first, then
spend calibration only where it says the oracle should be able to see. The procedure
does not replace calibration---it aims it.
